\subsection{Calendario de intervención y control de ventanas}
El escenario utiliza inicio el 1 de diciembre de 2026. Las habilitaciones externas tienen fecha requerida y responsable: recinto, equipos, red y telemetría recibidos para el 16 de marzo de 2027; licencias e identidad antes de su configuración; contratos contables y de pagos para enero de 2028. No se afirma que el CLIENTE haya ejecutado las obras. La instalación que afecte al varadero se prepara por anticipado y se ejecuta en la ventana de marzo posterior al día 15; la integración central y recepción de puntos de pantalán continúa en abril, que está abierto para el resto del recinto. H3 incluye cinco ambientes y la recepción de medición, con su reserva marina, antes del cierre de mayo.

Construcción E1 y componentes compartidos se concentran en agosto de 2027 con refuerzos dedicados; las vistas E1 ya no se dejan para septiembre. E2 aprueba modelos, contratos y esquemas en diciembre/enero, construye en enero--febrero y prueba desde el 16 de marzo de 2028. El trabajo remoto en DEV/QA no cambia producción ni habilita usuarios durante el congelamiento; toda intervención productiva, prueba destructiva o formación presencial respeta la ventana del sitio. Los ensayos usan manifiestos completos, dos ejecuciones por etapa y sus revisiones/subsanaciones. Una modificación que invalide evidencia obliga a repetirla sobre la versión recibida, a costo de Synaptix.

Las bases de IN1/IN4/IN5 abarcan 1--28 de mayo y la evaluación 1--28 de junio de 2028; en IN1 las solicitudes entran en los primeros 21 días y se siguen siete. La formación posterior a base se realiza 29--31 de mayo, con dos formadores y Calidad reservados. Las acciones de IN5 se aplican entre ambas ventanas, con autorización y evidencia. Los informes, revisión y subsanación se preparan en julio y las innovaciones se materializan en 21. IN3 conserva 120 pólizas distintas, orden intercalado y control de errores. IN2 tiene fecha concreta de invernada en junio de 2029 y controles de botadura antes de sus maniobras; sus actividades documentales no autorizan una intervención en varadero durante las campañas.

Los controles de DR se anticipan a noviembre, marzo después del día 15 y agosto, con nueve controles en 24/28/33/36/40/45/48/52/56, evitando exceder seis meses. Las pruebas ofensivas de operación se realizan en marzo después del día 15 de 2029, 2030 y 2031. Las jornadas de actualización y la incorporación estacional se preparan fuera del congelamiento; la autoformación no sustituye la evaluación. La salida se acompaña desde abril hasta julio de 2031, al menos noventa días reales, y la eliminación de origen depende de exportaciones, credenciales y acta de transferencia.

Proyecto registrará los catorce eventos anuales y sus bloqueos de uno a tres días, con fechas conocidas un año antes, y las fechas efectivas de las campañas de seis semanas. El escenario no inventa esas fechas. Las actividades de campo sólo se liberan tras comprobar zona, maniobra, evento, viento/marejada, acceso y autorización; ante aviso se cancelan y se consume reserva identificada. Si la fecha real de inicio o un evento anula la ventana, se recalculará la programación completa antes de fijar la línea base, sin trasladar los meses contractuales. La serie de medición comienza en E1 y se conserva antes de la temporada 2029; los residuos y la capacitación ambiental tienen recepción propia y no quedan demostrados sólo por instalar medidores.
